
\documentclass{article}
\usepackage{colortbl}
\usepackage{makecell}
\usepackage{multirow}
\usepackage{supertabular}

\begin{document}

\newcounter{utterance}

\centering \large Interaction Transcript for game `matchit\_ascii', experiment `similar\_grid\_2', episode 4 with nemotron{-}3{-}ultra.
\vspace{24pt}

{ \footnotesize  \setcounter{utterance}{1}
\setlength{\tabcolsep}{0pt}
\begin{supertabular}{c@{$\;$}|p{.15\linewidth}@{}p{.15\linewidth}p{.15\linewidth}p{.15\linewidth}p{.15\linewidth}p{.15\linewidth}}
    \# & $\;$A & \multicolumn{4}{c}{Game Master} & $\;\:$B\\
    \hline

    \theutterance \stepcounter{utterance}  
    & & \multicolumn{4}{p{0.6\linewidth}}{
        \cellcolor[rgb]{0.9,0.9,0.9}{
            \makecell[{{p{\linewidth}}}]{
                \texttt{\tiny{[P1$\langle$GM]}}
                \texttt{Du nimmst an einem gemeinsamen Ratespiel teil.} \\
\texttt{Das Ziel ist herauszufinden, ob dieses Raster und ein anderes Bild, das nur ich sehen kann, gleich sind.} \\
\texttt{Bitte beschreibe zuerst dein Raster. Dann gebe ich meine Beschreibung ab und wir können uns gegenseitig Fragen zu den Rastern stellen, um herauszufinden, ob sie gleich sind. Beginne jetzt deine kurze Rasterbeschreibung mit "BESCHREIBUNG:" gefolgt von der Beschreibung. Füge nichts anderes hinzu. Das ist dein Raster:} \\
\texttt{X X X X X} \\
\texttt{X ▢ ▢ ▢ X} \\
\texttt{X ▢ ▢ ▢ X} \\
\texttt{X ▢ ▢ ▢ X} \\
\texttt{X X X X X} \\
            }
        }
    }
    & \\ \\

    \theutterance \stepcounter{utterance}  
    & \multicolumn{4}{p{0.6\linewidth}}{
        \cellcolor[rgb]{0.9,0.9,0.9}{
            \makecell[{{p{\linewidth}}}]{
                \texttt{\tiny{[P1$\rangle$GM]}}
                \texttt{BESCHREIBUNG: Ein 5x5{-}Raster mit einem Rand aus X. Das innere 3x3{-}Feld besteht aus leeren Quadraten (▢).} \\
            }
        }
    }
    & & \\ \\

    \theutterance \stepcounter{utterance}  
    & & & \multicolumn{2}{p{0.3\linewidth}}{
        \cellcolor[rgb]{0.9,0.9,0.9}{
            \makecell[{{p{\linewidth}}}]{
                \texttt{\tiny{[GM$|$GM]}}
                \texttt{continue} \\
            }
        }
    }
    & & \\ \\

    \theutterance \stepcounter{utterance}  
    & & \multicolumn{4}{p{0.6\linewidth}}{
        \cellcolor[rgb]{0.9,0.9,0.9}{
            \makecell[{{p{\linewidth}}}]{
                \texttt{\tiny{[GM$\rangle$P2]}}
                \texttt{Du nimmst an einem gemeinsamen Ratespiel teil.} \\
\texttt{Das Ziel ist herauszufinden, ob dieses Raster und ein anderes Bild, das nur ich sehen kann, gleich sind.} \\
\texttt{Bitte beschreibe zuerst dein Raster. Dann gebe ich meine Beschreibung ab und wir können uns gegenseitig Fragen zu den Rastern stellen, um herauszufinden, ob sie gleich sind. Beginne jetzt deine kurze Rasterbeschreibung mit "BESCHREIBUNG:" gefolgt von der Beschreibung. Füge nichts anderes hinzu. Das ist dein Raster:} \\
\texttt{X ▢ X X X} \\
\texttt{X ▢ ▢ ▢ X} \\
\texttt{X ▢ ▢ ▢ X} \\
\texttt{X ▢ ▢ ▢ X} \\
\texttt{X ▢ X X X} \\
            }
        }
    }
    & \\ \\

    \theutterance \stepcounter{utterance}  
    & & & \multicolumn{4}{p{0.6\linewidth}}{
        \cellcolor[rgb]{0.9,0.9,0.9}{
            \makecell[{{p{\linewidth}}}]{
                \texttt{\tiny{[GM$\langle$P2]}}
                \texttt{BESCHREIBUNG: 5x5{-}Gitter. Rand oben und unten: X an den Ecken und in der Mitte, dazwischen ▢. Linke und rechte Spalte: komplett X. Inneres 3x3 (Zeilen 2{-}4, Spalten 2{-}4): komplett ▢.} \\
            }
        }
    }
     \\ \\

    \theutterance \stepcounter{utterance}  
    & & & \multicolumn{2}{p{0.3\linewidth}}{
        \cellcolor[rgb]{0.9,0.9,0.9}{
            \makecell[{{p{\linewidth}}}]{
                \texttt{\tiny{[GM$|$GM]}}
                \texttt{continue} \\
            }
        }
    }
    & & \\ \\

    \theutterance \stepcounter{utterance}  
    & & \multicolumn{4}{p{0.6\linewidth}}{
        \cellcolor[rgb]{0.9,0.9,0.9}{
            \makecell[{{p{\linewidth}}}]{
                \texttt{\tiny{[GM$\rangle$P2]}}
                \texttt{Das ist meineBESCHREIBUNG: Ein 5x5{-}Raster mit einem Rand aus X. Das innere 3x3{-}Feld besteht aus leeren Quadraten (▢).} \\
\texttt{Stelle jetzt eine Frage, um neue Aspekte meines Rasters herauszufinden, die sich möglicherweise von deinem Raster unterscheiden. Beginne mit "FRAGE:" und füge nichts anderes hinzu.} \\
            }
        }
    }
    & \\ \\

    \theutterance \stepcounter{utterance}  
    & & & \multicolumn{4}{p{0.6\linewidth}}{
        \cellcolor[rgb]{0.9,0.9,0.9}{
            \makecell[{{p{\linewidth}}}]{
                \texttt{\tiny{[GM$\langle$P2]}}
                \texttt{FRAGE: Sind in der ersten Reihe (oberer Rand) die Spalten 2 und 4 mit ▢ oder mit X gefüllt?} \\
            }
        }
    }
     \\ \\

    \theutterance \stepcounter{utterance}  
    & & & \multicolumn{2}{p{0.3\linewidth}}{
        \cellcolor[rgb]{0.9,0.9,0.9}{
            \makecell[{{p{\linewidth}}}]{
                \texttt{\tiny{[GM$|$GM]}}
                \texttt{continue} \\
            }
        }
    }
    & & \\ \\

    \theutterance \stepcounter{utterance}  
    & & \multicolumn{4}{p{0.6\linewidth}}{
        \cellcolor[rgb]{0.9,0.9,0.9}{
            \makecell[{{p{\linewidth}}}]{
                \texttt{\tiny{[P1$\langle$GM]}}
                \texttt{Das ist meineBESCHREIBUNG: 5x5{-}Gitter. Rand oben und unten: X an den Ecken und in der Mitte, dazwischen ▢. Linke und rechte Spalte: komplett X. Inneres 3x3 (Zeilen 2{-}4, Spalten 2{-}4): komplett ▢.} \\
\texttt{FRAGE: Sind in der ersten Reihe (oberer Rand) die Spalten 2 und 4 mit ▢ oder mit X gefüllt?Beginne deine Antwort mit "ANTWORT:" und füge nichts anderes hinzu.} \\
            }
        }
    }
    & \\ \\

    \theutterance \stepcounter{utterance}  
    & \multicolumn{4}{p{0.6\linewidth}}{
        \cellcolor[rgb]{0.9,0.9,0.9}{
            \makecell[{{p{\linewidth}}}]{
                \texttt{\tiny{[P1$\rangle$GM]}}
                \texttt{ANTWORT: In der ersten Reihe sind alle Spalten (1 bis 5) mit X gefüllt.} \\
            }
        }
    }
    & & \\ \\

    \theutterance \stepcounter{utterance}  
    & & & \multicolumn{2}{p{0.3\linewidth}}{
        \cellcolor[rgb]{0.9,0.9,0.9}{
            \makecell[{{p{\linewidth}}}]{
                \texttt{\tiny{[GM$|$GM]}}
                \texttt{continue} \\
            }
        }
    }
    & & \\ \\

    \theutterance \stepcounter{utterance}  
    & & \multicolumn{4}{p{0.6\linewidth}}{
        \cellcolor[rgb]{0.9,0.9,0.9}{
            \makecell[{{p{\linewidth}}}]{
                \texttt{\tiny{[P1$\langle$GM]}}
                \texttt{Stelle jetzt eine Frage, um neue Aspekte meines Rasters herauszufinden, die sich möglicherweise von deinem Raster unterscheiden. Beginne mit "FRAGE:" und füge nichts anderes hinzu.} \\
            }
        }
    }
    & \\ \\

    \theutterance \stepcounter{utterance}  
    & \multicolumn{4}{p{0.6\linewidth}}{
        \cellcolor[rgb]{0.9,0.9,0.9}{
            \makecell[{{p{\linewidth}}}]{
                \texttt{\tiny{[P1$\rangle$GM]}}
                \texttt{FRAGE: Sind in der fünften Reihe (unterer Rand) die Spalten 2 und 4 mit ▢ oder mit X gefüllt?} \\
            }
        }
    }
    & & \\ \\

    \theutterance \stepcounter{utterance}  
    & & & \multicolumn{2}{p{0.3\linewidth}}{
        \cellcolor[rgb]{0.9,0.9,0.9}{
            \makecell[{{p{\linewidth}}}]{
                \texttt{\tiny{[GM$|$GM]}}
                \texttt{continue} \\
            }
        }
    }
    & & \\ \\

    \theutterance \stepcounter{utterance}  
    & & \multicolumn{4}{p{0.6\linewidth}}{
        \cellcolor[rgb]{0.9,0.9,0.9}{
            \makecell[{{p{\linewidth}}}]{
                \texttt{\tiny{[GM$\rangle$P2]}}
                \texttt{ANTWORT: In der ersten Reihe sind alle Spalten (1 bis 5) mit X gefüllt.} \\
\texttt{FRAGE: Sind in der fünften Reihe (unterer Rand) die Spalten 2 und 4 mit ▢ oder mit X gefüllt?Beginne deine Antwort mit "ANTWORT:" und füge nichts anderes hinzu.} \\
            }
        }
    }
    & \\ \\

    \theutterance \stepcounter{utterance}  
    & & & \multicolumn{4}{p{0.6\linewidth}}{
        \cellcolor[rgb]{0.9,0.9,0.9}{
            \makecell[{{p{\linewidth}}}]{
                \texttt{\tiny{[GM$\langle$P2]}}
                \texttt{ANTWORT: In der fünften Reihe sind Spalte 2 und 4 mit X gefüllt.} \\
\texttt{FRAGE: Wie sind die Ecken des Gitters (Positionen (1,1), (1,5), (5,1), (5,5)) belegt?} \\
            }
        }
    }
     \\ \\

    \theutterance \stepcounter{utterance}  
    & & & \multicolumn{2}{p{0.3\linewidth}}{
        \cellcolor[rgb]{0.9,0.9,0.9}{
            \makecell[{{p{\linewidth}}}]{
                \texttt{\tiny{[GM$|$GM]}}
                \texttt{continue} \\
            }
        }
    }
    & & \\ \\

    \theutterance \stepcounter{utterance}  
    & & \multicolumn{4}{p{0.6\linewidth}}{
        \cellcolor[rgb]{0.9,0.9,0.9}{
            \makecell[{{p{\linewidth}}}]{
                \texttt{\tiny{[GM$\rangle$P2]}}
                \texttt{Stelle jetzt eine Frage, um neue Aspekte meines Rasters herauszufinden, die sich möglicherweise von deinem Raster unterscheiden. Beginne mit "FRAGE:" und füge nichts anderes hinzu.} \\
            }
        }
    }
    & \\ \\

    \theutterance \stepcounter{utterance}  
    & & & \multicolumn{4}{p{0.6\linewidth}}{
        \cellcolor[rgb]{0.9,0.9,0.9}{
            \makecell[{{p{\linewidth}}}]{
                \texttt{\tiny{[GM$\langle$P2]}}
                \texttt{FRAGE: Gibt es in deinem Gitter irgendwo ein X, das nicht zum äußeren Rand gehört (also in Zeilen 2{-}4 und Spalten 2{-}4)?} \\
            }
        }
    }
     \\ \\

    \theutterance \stepcounter{utterance}  
    & & & \multicolumn{2}{p{0.3\linewidth}}{
        \cellcolor[rgb]{0.9,0.9,0.9}{
            \makecell[{{p{\linewidth}}}]{
                \texttt{\tiny{[GM$|$GM]}}
                \texttt{continue} \\
            }
        }
    }
    & & \\ \\

    \theutterance \stepcounter{utterance}  
    & & \multicolumn{4}{p{0.6\linewidth}}{
        \cellcolor[rgb]{0.9,0.9,0.9}{
            \makecell[{{p{\linewidth}}}]{
                \texttt{\tiny{[P1$\langle$GM]}}
                \texttt{ANTWORT: In der fünften Reihe sind Spalte 2 und 4 mit X gefüllt.} \\
\texttt{FRAGE: Wie sind die Ecken des Gitters (Positionen (1,1), (1,5), (5,1), (5,5)) belegt?} \\
\texttt{FRAGE: Gibt es in deinem Gitter irgendwo ein X, das nicht zum äußeren Rand gehört (also in Zeilen 2{-}4 und Spalten 2{-}4)?Beginne deine Antwort mit "ANTWORT:" und füge nichts anderes hinzu.} \\
            }
        }
    }
    & \\ \\

    \theutterance \stepcounter{utterance}  
    & \multicolumn{4}{p{0.6\linewidth}}{
        \cellcolor[rgb]{0.9,0.9,0.9}{
            \makecell[{{p{\linewidth}}}]{
                \texttt{\tiny{[P1$\rangle$GM]}}
                \texttt{ANTWORT: Alle vier Ecken sind mit X belegt. Im inneren 3x3{-}Bereich (Zeilen 2{-}4, Spalten 2{-}4) gibt es keine X, dort sind alle Felder ▢.} \\
            }
        }
    }
    & & \\ \\

    \theutterance \stepcounter{utterance}  
    & & & \multicolumn{2}{p{0.3\linewidth}}{
        \cellcolor[rgb]{0.9,0.9,0.9}{
            \makecell[{{p{\linewidth}}}]{
                \texttt{\tiny{[GM$|$GM]}}
                \texttt{continue} \\
            }
        }
    }
    & & \\ \\

    \theutterance \stepcounter{utterance}  
    & & \multicolumn{4}{p{0.6\linewidth}}{
        \cellcolor[rgb]{0.9,0.9,0.9}{
            \makecell[{{p{\linewidth}}}]{
                \texttt{\tiny{[P1$\langle$GM]}}
                \texttt{Stelle jetzt eine Frage, um neue Aspekte meines Rasters herauszufinden, die sich möglicherweise von deinem Raster unterscheiden. Beginne mit "FRAGE:" und füge nichts anderes hinzu.} \\
            }
        }
    }
    & \\ \\

    \theutterance \stepcounter{utterance}  
    & \multicolumn{4}{p{0.6\linewidth}}{
        \cellcolor[rgb]{0.9,0.9,0.9}{
            \makecell[{{p{\linewidth}}}]{
                \texttt{\tiny{[P1$\rangle$GM]}}
                \texttt{FRAGE: Wie sind die Positionen (1,3) und (5,3) (Mitte der oberen und unteren Reihe) belegt?} \\
            }
        }
    }
    & & \\ \\

    \theutterance \stepcounter{utterance}  
    & & & \multicolumn{2}{p{0.3\linewidth}}{
        \cellcolor[rgb]{0.9,0.9,0.9}{
            \makecell[{{p{\linewidth}}}]{
                \texttt{\tiny{[GM$|$GM]}}
                \texttt{continue} \\
            }
        }
    }
    & & \\ \\

    \theutterance \stepcounter{utterance}  
    & & \multicolumn{4}{p{0.6\linewidth}}{
        \cellcolor[rgb]{0.9,0.9,0.9}{
            \makecell[{{p{\linewidth}}}]{
                \texttt{\tiny{[GM$\rangle$P2]}}
                \texttt{ANTWORT: Alle vier Ecken sind mit X belegt. Im inneren 3x3{-}Bereich (Zeilen 2{-}4, Spalten 2{-}4) gibt es keine X, dort sind alle Felder ▢.} \\
\texttt{FRAGE: Wie sind die Positionen (1,3) und (5,3) (Mitte der oberen und unteren Reihe) belegt?Beginne deine Antwort mit "ANTWORT:" und füge nichts anderes hinzu.} \\
            }
        }
    }
    & \\ \\

    \theutterance \stepcounter{utterance}  
    & & & \multicolumn{4}{p{0.6\linewidth}}{
        \cellcolor[rgb]{0.9,0.9,0.9}{
            \makecell[{{p{\linewidth}}}]{
                \texttt{\tiny{[GM$\langle$P2]}}
                \texttt{ANTWORT: Position (1,3) und (5,3) sind mit X belegt.} \\
\texttt{FRAGE: Sind die Positionen (3,1) und (3,5) (Mitte der linken und rechten Spalte) mit X oder ▢ belegt?} \\
            }
        }
    }
     \\ \\

    \theutterance \stepcounter{utterance}  
    & & & \multicolumn{2}{p{0.3\linewidth}}{
        \cellcolor[rgb]{0.9,0.9,0.9}{
            \makecell[{{p{\linewidth}}}]{
                \texttt{\tiny{[GM$|$GM]}}
                \texttt{continue} \\
            }
        }
    }
    & & \\ \\

    \theutterance \stepcounter{utterance}  
    & & \multicolumn{4}{p{0.6\linewidth}}{
        \cellcolor[rgb]{0.9,0.9,0.9}{
            \makecell[{{p{\linewidth}}}]{
                \texttt{\tiny{[GM$\rangle$P2]}}
                \texttt{Stelle jetzt eine Frage, um neue Aspekte meines Rasters herauszufinden, die sich möglicherweise von deinem Raster unterscheiden. Beginne mit "FRAGE:" und füge nichts anderes hinzu.} \\
            }
        }
    }
    & \\ \\

    \theutterance \stepcounter{utterance}  
    & & & \multicolumn{4}{p{0.6\linewidth}}{
        \cellcolor[rgb]{0.9,0.9,0.9}{
            \makecell[{{p{\linewidth}}}]{
                \texttt{\tiny{[GM$\langle$P2]}}
                \texttt{FRAGE: Sind alle vier Seitenränder (obere Reihe, untere Reihe, linke Spalte, rechte Spalte) vollständig und lückenlos mit X gefüllt?} \\
            }
        }
    }
     \\ \\

    \theutterance \stepcounter{utterance}  
    & & & \multicolumn{2}{p{0.3\linewidth}}{
        \cellcolor[rgb]{0.9,0.9,0.9}{
            \makecell[{{p{\linewidth}}}]{
                \texttt{\tiny{[GM$|$GM]}}
                \texttt{continue} \\
            }
        }
    }
    & & \\ \\

    \theutterance \stepcounter{utterance}  
    & & \multicolumn{4}{p{0.6\linewidth}}{
        \cellcolor[rgb]{0.9,0.9,0.9}{
            \makecell[{{p{\linewidth}}}]{
                \texttt{\tiny{[P1$\langle$GM]}}
                \texttt{ANTWORT: Position (1,3) und (5,3) sind mit X belegt.} \\
\texttt{FRAGE: Sind die Positionen (3,1) und (3,5) (Mitte der linken und rechten Spalte) mit X oder ▢ belegt?} \\
\texttt{FRAGE: Sind alle vier Seitenränder (obere Reihe, untere Reihe, linke Spalte, rechte Spalte) vollständig und lückenlos mit X gefüllt?Beginne deine Antwort mit "ANTWORT:" und füge nichts anderes hinzu.} \\
            }
        }
    }
    & \\ \\

    \theutterance \stepcounter{utterance}  
    & \multicolumn{4}{p{0.6\linewidth}}{
        \cellcolor[rgb]{0.9,0.9,0.9}{
            \makecell[{{p{\linewidth}}}]{
                \texttt{\tiny{[P1$\rangle$GM]}}
                \texttt{ANTWORT: Ja, Positionen (3,1) und (3,5) sind mit X belegt. Ja, alle vier Seitenränder sind vollständig und lückenlos mit X gefüllt.} \\
            }
        }
    }
    & & \\ \\

    \theutterance \stepcounter{utterance}  
    & & & \multicolumn{2}{p{0.3\linewidth}}{
        \cellcolor[rgb]{0.9,0.9,0.9}{
            \makecell[{{p{\linewidth}}}]{
                \texttt{\tiny{[GM$|$GM]}}
                \texttt{continue} \\
            }
        }
    }
    & & \\ \\

    \theutterance \stepcounter{utterance}  
    & & \multicolumn{4}{p{0.6\linewidth}}{
        \cellcolor[rgb]{0.9,0.9,0.9}{
            \makecell[{{p{\linewidth}}}]{
                \texttt{\tiny{[P1$\langle$GM]}}
                \texttt{Stelle jetzt eine Frage, um neue Aspekte meines Rasters herauszufinden, die sich möglicherweise von deinem Raster unterscheiden. Beginne mit "FRAGE:" und füge nichts anderes hinzu.} \\
            }
        }
    }
    & \\ \\

    \theutterance \stepcounter{utterance}  
    & \multicolumn{4}{p{0.6\linewidth}}{
        \cellcolor[rgb]{0.9,0.9,0.9}{
            \makecell[{{p{\linewidth}}}]{
                \texttt{\tiny{[P1$\rangle$GM]}}
                \texttt{FRAGE: Gibt es in deinem Gitter Symbole außer X und ▢?} \\
            }
        }
    }
    & & \\ \\

    \theutterance \stepcounter{utterance}  
    & & & \multicolumn{2}{p{0.3\linewidth}}{
        \cellcolor[rgb]{0.9,0.9,0.9}{
            \makecell[{{p{\linewidth}}}]{
                \texttt{\tiny{[GM$|$GM]}}
                \texttt{continue} \\
            }
        }
    }
    & & \\ \\

    \theutterance \stepcounter{utterance}  
    & & \multicolumn{4}{p{0.6\linewidth}}{
        \cellcolor[rgb]{0.9,0.9,0.9}{
            \makecell[{{p{\linewidth}}}]{
                \texttt{\tiny{[GM$\rangle$P2]}}
                \texttt{ANTWORT: Ja, Positionen (3,1) und (3,5) sind mit X belegt. Ja, alle vier Seitenränder sind vollständig und lückenlos mit X gefüllt.} \\
\texttt{FRAGE: Gibt es in deinem Gitter Symbole außer X und ▢?Beginne deine Antwort mit "ANTWORT:" und füge nichts anderes hinzu.} \\
            }
        }
    }
    & \\ \\

    \theutterance \stepcounter{utterance}  
    & & & \multicolumn{4}{p{0.6\linewidth}}{
        \cellcolor[rgb]{0.9,0.9,0.9}{
            \makecell[{{p{\linewidth}}}]{
                \texttt{\tiny{[GM$\langle$P2]}}
                \texttt{ANTWORT: Nein, es gibt nur X und ▢.} \\
\texttt{FRAGE: Stimmt die Beschreibung "Rand komplett X, Inneres 3x3 komplett ▢" exakt mit deinem Gitter überein?} \\
            }
        }
    }
     \\ \\

    \theutterance \stepcounter{utterance}  
    & & & \multicolumn{2}{p{0.3\linewidth}}{
        \cellcolor[rgb]{0.9,0.9,0.9}{
            \makecell[{{p{\linewidth}}}]{
                \texttt{\tiny{[GM$|$GM]}}
                \texttt{continue} \\
            }
        }
    }
    & & \\ \\

    \theutterance \stepcounter{utterance}  
    & & \multicolumn{4}{p{0.6\linewidth}}{
        \cellcolor[rgb]{0.9,0.9,0.9}{
            \makecell[{{p{\linewidth}}}]{
                \texttt{\tiny{[GM$\rangle$P2]}}
                \texttt{Triff jetzt eine Entscheidung. Was denkst du: Sind dein Raster und das andere beschriebene Raster identisch? Schreibe "ENTSCHEIDUNG: gleiches Raster", wenn du denkst, dass es dasselbe Raster ist, oder "ENTSCHEIDUNG: verschiedene Raster", wenn du denkst, dass es unterschiedliche Raster sind. Füge nichts anderes hinzu.} \\
            }
        }
    }
    & \\ \\

    \theutterance \stepcounter{utterance}  
    & & & \multicolumn{4}{p{0.6\linewidth}}{
        \cellcolor[rgb]{0.9,0.9,0.9}{
            \makecell[{{p{\linewidth}}}]{
                \texttt{\tiny{[GM$\langle$P2]}}
                \texttt{ENTSCHEIDUNG: verschiedene Raster} \\
            }
        }
    }
     \\ \\

    \theutterance \stepcounter{utterance}  
    & & & \multicolumn{2}{p{0.3\linewidth}}{
        \cellcolor[rgb]{0.9,0.9,0.9}{
            \makecell[{{p{\linewidth}}}]{
                \texttt{\tiny{[GM$|$GM]}}
                \texttt{continue} \\
            }
        }
    }
    & & \\ \\

    \theutterance \stepcounter{utterance}  
    & & & \multicolumn{2}{p{0.3\linewidth}}{
        \cellcolor[rgb]{0.9,0.9,0.9}{
            \makecell[{{p{\linewidth}}}]{
                \texttt{\tiny{[GM$|$GM]}}
                \texttt{success} \\
            }
        }
    }
    & & \\ \\

    \theutterance \stepcounter{utterance}  
    & & \multicolumn{4}{p{0.6\linewidth}}{
        \cellcolor[rgb]{0.9,0.9,0.9}{
            \makecell[{{p{\linewidth}}}]{
                \texttt{\tiny{[P1$\langle$GM]}}
                \texttt{ANTWORT: Nein, es gibt nur X und ▢.} \\
\texttt{FRAGE: Stimmt die Beschreibung "Rand komplett X, Inneres 3x3 komplett ▢" exakt mit deinem Gitter überein?} \\
\texttt{Triff jetzt eine Entscheidung. Was denkst du: Sind dein Raster und das andere beschriebene Raster identisch? Schreibe "ENTSCHEIDUNG: gleiches Raster", wenn du denkst, dass es dasselbe Raster ist, oder "ENTSCHEIDUNG: verschiedene Raster", wenn du denkst, dass es unterschiedliche Raster sind. Füge nichts anderes hinzu.} \\
            }
        }
    }
    & \\ \\

    \theutterance \stepcounter{utterance}  
    & \multicolumn{4}{p{0.6\linewidth}}{
        \cellcolor[rgb]{0.9,0.9,0.9}{
            \makecell[{{p{\linewidth}}}]{
                \texttt{\tiny{[P1$\rangle$GM]}}
                \texttt{ENTSCHEIDUNG: gleiches Raster} \\
            }
        }
    }
    & & \\ \\

    \theutterance \stepcounter{utterance}  
    & & & \multicolumn{2}{p{0.3\linewidth}}{
        \cellcolor[rgb]{0.9,0.9,0.9}{
            \makecell[{{p{\linewidth}}}]{
                \texttt{\tiny{[GM$|$GM]}}
                \texttt{continue} \\
            }
        }
    }
    & & \\ \\

    \theutterance \stepcounter{utterance}  
    & & & \multicolumn{2}{p{0.3\linewidth}}{
        \cellcolor[rgb]{0.9,0.9,0.9}{
            \makecell[{{p{\linewidth}}}]{
                \texttt{\tiny{[GM$|$GM]}}
                \texttt{loss} \\
            }
        }
    }
    & & \\ \\

\end{supertabular}
}

\end{document}
